% ch5_d 第5章 §5.4尾–§5.5.4 (书页 226–238 / p241–p253)
%JOIN%
两个电荷由于来自其他电子的屏蔽，并不感受到长程相互作用。因此，人们可能会怀疑 Hartree–Fock 关于发散费米速度这一结果的正确性；尤其是在认识到 Hartree–Fock 近似只包含直接相互作用、不含任何屏蔽效应之后，更会如此。

%FIG{5.13}: {(a) 相互作用产生粒子–空穴激发；(b) 粒子–空穴激发修正相互作用。}

%FIG{5.14}: {无规相位近似下的有效势。}

为了超越 Hartree–Fock 近似并把屏蔽效应包括进来，我们必须包括这样的效应：相互作用产生粒子–空穴激发（图 5.13(a)），而粒子–空穴激发又转而修正势（图 5.13(b)）。在无规相位近似（RPA）中，我们把图 5.14 中的图包括进来，如下计算屏蔽后的有效势：
\begin{equation*}
-\mathrm{i}V_{\pmb{q}\nu}^{\mathrm{eff}}=-\mathrm{i}V_{\pmb{q}\nu}\left(1+(-\mathrm{i}V_{\pmb{q}\nu})(\mathrm{i}P^{00}_{\pmb{q}\nu})+(-\mathrm{i}V_{\pmb{q}\nu})^{2}(\mathrm{i}P^{00}_{\pmb{q}\nu})^{2}+\cdots\right)
\end{equation*}
我们发现
\begin{equation*}
V_{\pmb{q}\nu}^{\mathrm{eff}}=\frac{V_{\pmb{q}\nu}}{1-V_{\pmb{q}\nu}P^{00}_{\pmb{q}\nu}}
\end{equation*}
在 $\nu\ll v_{F}q$ 极限下，$P^{00}_{\pmb{q}\nu}$ 等于压缩率的负值，即 $P^{00}_{\pmb{q}\nu}=-N(E_{F})$。我们有
\begin{equation*}
V_{\pmb{q}\nu}^{\mathrm{eff}}=\frac{V_{\pmb{q}\nu}}{1+V_{\pmb{q}\nu}N(E_{F})}
\end{equation*}
而对库仑相互作用，我们有
\begin{equation*}
V^{\mathrm{eff}}_{\pmb{q}}=\frac{4\pi e^{2}}{\pmb{q}^{2}+4\pi e^{2}N(E_{F})}
\end{equation*}

%FIG{5.15}: {准粒子的衰变。}

显然，在 $\nu\ll v_{F}q$ 极限下，有效势在 $\pmb{q}\to 0$ 时并不奇异。这样的有效势代表一个屏蔽了的短程相互作用。同样的结果也可以通过 Thomas–Fermi 模型得到。

RPA 计算与屏蔽之间的关系可以按下面的方式看得更清楚。RPA 计算中的第一项（见图 5.14）代表两个不同位置处的两个密度（在图 5.14 中用两段短竖线表示）之间的直接（即未屏蔽的）相互作用。RPA 计算中的第二项代表这样的效应：一处的密度通过相互作用 $V_{\pmb{q}\nu}$ 诱导出一个粒子–空穴激发。该粒子–空穴激发代表一份形变了的密度，它转而与另一处的密度发生相互作用。通过这种方式，第二项（以及其他更高阶的项）修正或屏蔽了两个位置处的密度之间的相互作用。

在标准实践中，我们把裸势 $V_{\pmb{q}\nu}$ 换成屏蔽势 $V^{\mathrm{eff}}_{\pmb{q}\nu}=V_{\pmb{q}\nu}/(1+V_{\pmb{q}\nu}N(E_{F}))$，然后用这个有效势去计算自能与费米液体函数。在这一近似中，我们忽略了 $P^{00}$ 的频率依赖。一个瞬时的裸相互作用（不依赖频率的相互作用）被一个瞬时的有效相互作用所近似。如果把 $P^{00}$ 的频率依赖包括进来，那么一个瞬时的裸相互作用可以导致一个依赖频率的有效相互作用。这种频率依赖代表的正是推迟效应（retardation effects）。

\subsection{Landau 费米液体理论的论证}

\begin{keypoints}
\item 相互作用使准粒子发生衰变，但衰变率远小于准粒子的能量。Landau 费米液体理论尚待验证。
\end{keypoints}

一个电子–电子相互作用可以使动量为 $\pmb{k}$ 的电子衰变为一个动量为 $\pmb{q}$ 的电子加上一个动量为 $\pmb{k}-\pmb{q}$ 的粒子–空穴激发（见图 5.15）。粒子–空穴激发中的粒子与空穴各自可以具有 0 到 $\xi_{\pmb{k}}$ 之间的能量。$\pmb{q}$ 处粒子的能量会相应地调整以维持能量守恒。因此，可用衰变通道的数目正比于初始电子能量 $\xi_{\pmb{k}}$ 的平方。于是衰变率也正比于 $\xi^{2}_{\pmb{k}}$。由于衰变振幅正比于 $V_{\pmb{k}-\pmb{q}}$，我们求得衰变率为
\begin{equation*}
\Gamma_{\pmb{k}}\propto|V_{\pmb{k}}|^{2}\xi^{2}_{\pmb{k}}
\end{equation*}

%FIG{5.16}: {对自能的 $V_{\pmb{q}}^{2}$ 量级的一项贡献。该图由两半构成，每一半代表图 5.15 中的衰变过程。上图给出衰变率，而图 5.15 中的图给出衰变振幅。}

准粒子的衰变可以由自能的虚部来描述。当 $\Sigma_{\pmb{k}\omega}$ 具有虚部时，电子格林函数 $1/(\omega-\xi_{\pmb{k}}-\Sigma_{\pmb{k}\omega})$ 的极点就偏离实 $\omega$ 轴。这导致实时格林函数中的衰变：
\begin{equation*}
G(t,\pmb{k})=\int\frac{\mathrm{d}\omega}{2\pi}\,\frac{1}{\omega-\xi_{\pmb{k}}-\Sigma_{\pmb{k}\omega}}\,\mathrm{e}^{-\mathrm{i}\omega t}\sim\mathrm{e}^{-\mathrm{i}\xi^{*}_{\pmb{k}}t}\,\mathrm{e}^{-|\mathrm{Im}\Sigma_{\pmb{k}\xi^{*}_{\pmb{k}}}|t}
\end{equation*}
其中 $\xi^{*}_{\pmb{k}}=\xi_{\pmb{k}}+\mathrm{Re}\Sigma_{\pmb{k}\xi^{*}_{\pmb{k}}}$。我们看到自能的虚部由下式给出：
\begin{equation*}
|\mathrm{Im}\Sigma_{\pmb{k}\xi^{*}_{\pmb{k}}}|=\Gamma_{\pmb{k}}
\end{equation*}
我们想指出，这样的自能可以由图 5.16 中的费曼图得到。

由于衰变率 $\Gamma=\mathrm{Im}\Sigma_{\pmb{k}\xi^{*}_{\pmb{k}}}\propto(\xi^{*}_{\pmb{k}})^{2}$ 在低能极限下远小于准粒子能量 $\xi^{*}_{\pmb{k}}$，我们可以忽略准粒子的衰变。在这一极限下，准粒子激发可以被看作本征态。这就是 Landau 费米液体理论的基础。我们看到 $\mathrm{Im}G_{\pmb{k}\omega}$ 在 $\xi^{*}_{\pmb{k}}$ 处几乎是一个 $\delta$ 函数，峰的宽度仅正比于 $(\xi^{*}_{\pmb{k}})^{2}$。

在费米面附近我们可以使用展开
\begin{equation*}
\mathrm{Re}\Sigma_{\pmb{k}\omega}=\big(1-\frac{1}{Z}\big)\omega+\big(\frac{v^{*}_{F}}{Z}-v_{F}\big)(k-k_{F})
\end{equation*}
并把格林函数改写为
\begin{equation*}
G_{\pmb{k}\omega}=\frac{1}{\omega-v_{F}(k-k_{F})-\mathrm{Re}\Sigma_{\pmb{k}\omega}}=\frac{Z}{\omega-v^{*}_{F}(k-k_{F})}
\end{equation*}
这里 $Z$ 是 5.3.1 节末尾引入的交叠因子。这样，我们就在一个微扰计算之内论证了 Landau 费米液体理论。

\section{对称破缺相与自旋密度波态}

\begin{keypoints}
\item 费米液体可以有许多失稳。相互作用可以导致对称破缺相。
\end{keypoints}

相互作用费米子系统可以有许多不同的态，这些态造就了性质如此繁多的材料，例如超导体、磁体、电荷密度波态等等。材料的所有这些不同的态都可以理解为相互作用诱导的某些对称破缺态。本节中我们将只讨论其中一种对称破缺态——自旋密度波态。这里讨论的图像、方法与结果可以很容易地推广到其他对称破缺态，例如超导态。

\subsection{线性响应与失稳}

\begin{keypoints}
\item 一个具有如此众多无能隙激发的自由费米子系统，拥有许多潜在的失稳。费米子之间的相互作用可以把潜在的失稳变成真实的失稳。不同的相互作用诱导不同的失稳。
\end{keypoints}

考虑一个具有在位相互作用的格点电子系统：
\begin{align*}
H=&-\sum_{\langle ij\rangle,\alpha}t(c^{\dagger}_{\alpha j}c_{\alpha i}+\text{h.c.})-\mu N+U\sum_{i}n_{i}^{2}\\
=&-\sum_{\langle ij\rangle,\alpha}t(c^{\dagger}_{\alpha j}c_{\alpha i}+\text{h.c.})-\mu' N+U\sum_{i}c^{\dagger}_{\beta i}c^{\dagger}_{\alpha i}c_{\alpha i}c_{\beta i}\tag{5.5.1}
\end{align*}
其中 $n_{i}=\sum_{\alpha}c^{\dagger}_{\alpha i}c_{\alpha i}$ 是格点 $i$ 上的电子数，$N=\sum_{i}n_{i}$ 是电子总数。上述模型称为 Hubbard 模型。

下面我们将集中讨论二维的 Hubbard 模型。在动量空间中，二维 Hubbard 模型可以改写为
\begin{equation*}
H=\sum_{\pmb{k}}\xi_{\pmb{k}}c^{\dagger}_{\alpha\pmb{k}}c_{\alpha\pmb{k}}+U\frac{1}{N_{\mathrm{site}}}\sum_{\pmb{k},\pmb{k}',\pmb{q}}c^{\dagger}_{\alpha\pmb{k}}c^{\dagger}_{\beta\pmb{k}'}c_{\beta,\pmb{k}'+\pmb{q}}c_{\alpha,\pmb{k}-\pmb{q}}
\end{equation*}
其中
\begin{equation*}
\xi_{\pmb{k}}=-2t\left(\cos(k_{x})+\cos(k_{y})\right)-\mu'
\end{equation*}
$N_{\mathrm{site}}$ 是格点数目。我们注意到，当 $\mu'=0$ 时，费米海是半满的（$\langle n_{i}\rangle=1$），而且费米面是一个正方形（见图 5.17），它满足嵌套条件，嵌套波矢为 $\pmb{Q}=(\pi,\pi)$。由式 (4.3.3)，我们看到自由费米子在嵌套波矢 $\pmb{Q}$ 处的压缩率与自旋磁化率在低温下发散。由于态密度 $N(E_{F})$（这里 $N(E_{F})$ 指每格点每单位能量的态数）也具有对数发散，我们有

%FIG{5.17}: {Hubbard 模型中的电子色散 $\xi_{\pmb{k}}$。}

%FIG{5.18}: {RPA 下的压缩率与自旋磁化率。}

\begin{equation*}
\chi_{0}(\pmb{Q})\sim\frac{1}{E_{F}}\big(\ln\frac{E_{F}}{T}\big)^{2}\tag{5.5.2}
\end{equation*}
$T=0$ 处发散的 $\chi_{0}(\pmb{Q})$ 意味着系统正濒于产生密度序或自旋序。下面我们将证明，对于相互作用系统，当 $T$ 趋于一个有限的临界温度 $T_{c}$ 时，磁化率 $\chi(\pmb{Q})$ 发散。在 $T_{c}$ 以下，系统自发地产生电荷密度波（CDW）态或自旋密度波（SDW）态。

对于相互作用电子，压缩率与自旋磁化率会收到来自相互作用的修正。在 RPA 中，我们只考虑由图 5.18 中的图所描述的修正。单圈气泡图由如下关联给出：
\begin{equation*}
\left\langle \rho_{\alpha;\pmb{q}\nu}\rho_{\beta;-\pmb{q},-\nu}\right\rangle=\mathrm{i}\Pi^{00}_{\alpha\beta;\pmb{q}\nu}=\mathrm{i}\Pi^{00}_{\pmb{q}\nu}\delta_{\alpha\beta}
\end{equation*}

%FIG{5.19}: {对自旋与密度关联有贡献的梯形图。}

而相互作用线由下式给出：
\begin{equation*}
-\mathrm{i}V_{\alpha\beta;\pmb{q}\nu}=-\mathrm{i}UC_{\alpha\beta}
\end{equation*}
其中 $\rho_{\alpha}$ 是自旋 $\alpha$ 电子的密度，$\pmb{C}=(C_{\alpha\beta})$ 是一个所有元素都等于 1 的 $2\times 2$ 矩阵。我们发现在 RPA 下，
\begin{align*}
\mathrm{i}\Pi^{\mathrm{RPA}}_{\alpha\beta;\pmb{q}\nu}&=\mathrm{i}\Pi^{00}_{\alpha\beta;\pmb{q}\nu}+\sum_{\gamma,\lambda}\mathrm{i}\Pi^{00}_{\alpha\gamma;\pmb{q}\nu}(-\mathrm{i}V_{\gamma\lambda;\pmb{q}\nu})\mathrm{i}\Pi^{00}_{\lambda\beta;\pmb{q}\nu}+\cdots\cdots\\
&=\sum_{\gamma}\mathrm{i}\Pi^{00}_{\pmb{q}\nu}\big(1-U\Pi^{00}_{\pmb{q}\nu}\pmb{C}\big)^{-1}_{\alpha\beta}=\mathrm{i}\Pi^{00}_{\pmb{q}\nu}\Big(\delta_{\alpha\beta}+\frac{U\Pi^{00}_{\pmb{q}\nu}}{1-2U\Pi^{00}_{\pmb{q}\nu}}C_{\alpha\beta}\Big)
\end{align*}
其中我们用到了 $(a+b\pmb{C})^{-1}=a^{-1}-b(a^{2}+2ab)^{-1}\pmb{C}$。于是，密度响应函数由下式给出：
\begin{align*}
\mathrm{i}\Pi^{\mathrm{RPA}}_{\mathrm{den};\pmb{q}\nu}&=\left\langle (\rho_{\uparrow;\pmb{q}\nu}+\rho_{\downarrow;\pmb{q}\nu})(\rho_{\uparrow;-\pmb{q},-\nu}+\rho_{\downarrow;-\pmb{q},-\nu})\right\rangle\\
&=\mathrm{i}\Pi^{\mathrm{RPA}}_{\uparrow\uparrow;\pmb{q}\nu}+\mathrm{i}\Pi^{\mathrm{RPA}}_{\downarrow\downarrow;\pmb{q}\nu}+2\,\mathrm{i}\Pi^{\mathrm{RPA}}_{\uparrow\downarrow;\pmb{q}\nu}=2\,\mathrm{i}\Pi^{00}_{\pmb{q}\nu}\frac{1}{1-2U\Pi^{00}_{\pmb{q}\nu}}
\end{align*}
而自旋响应函数由下式给出：
\begin{align*}
\mathrm{i}\Pi^{\mathrm{RPA}}_{\mathrm{spin};\pmb{q}\nu}&=\left\langle \frac{1}{2}(\rho_{\uparrow;\pmb{q}\nu}-\rho_{\downarrow;\pmb{q}\nu})\frac{1}{2}(\rho_{\uparrow;-\pmb{q},-\nu}-\rho_{\downarrow;-\pmb{q},-\nu})\right\rangle\\
&=\frac{1}{4}\,\mathrm{i}\Pi^{\mathrm{RPA}}_{\uparrow\uparrow;\pmb{q}\nu}+\frac{1}{4}\,\mathrm{i}\Pi^{\mathrm{RPA}}_{\downarrow\downarrow;\pmb{q}\nu}-\frac{1}{2}\,\mathrm{i}\Pi^{\mathrm{RPA}}_{\uparrow\downarrow;\pmb{q}\nu}=\frac{1}{2}\,\mathrm{i}\Pi^{00}_{\pmb{q}\nu}
\end{align*}

然而，上述结果并不十分完全。对于在位相互作用，图 5.19 中的梯形图与图 5.18 中的图有类似的贡献。把梯形图与气泡图的混合图包括进来，我们得到
\begin{align*}
\mathrm{i}\Pi^{\mathrm{RPA}^{+}}_{\alpha\beta;\pmb{q}\nu}&=\mathrm{i}\Pi^{00}_{\pmb{q}\nu}\delta_{\alpha\beta}+\sum_{\gamma\lambda}\left(\mathrm{i}\Pi^{00}_{\pmb{q}\nu}\delta_{\alpha\gamma}(-\mathrm{i}UC_{\gamma\lambda})\mathrm{i}\Pi^{00}_{\pmb{q}\nu}\delta_{\lambda\beta}\right)\\
&\quad+\sum_{\gamma\lambda}\left((-)\mathrm{i}\Pi^{00}_{\pmb{q}\nu}\delta_{\alpha\gamma}(-\mathrm{i}U\delta_{\gamma\lambda})\mathrm{i}\Pi^{00}_{\pmb{q}\nu}\delta_{\lambda\beta}\right)+\cdots\\
&=\sum_{\gamma}\mathrm{i}\Pi^{00}_{\pmb{q}\nu}\left(1-U\Pi^{00}_{\pmb{q}\nu}\pmb{C}+U\Pi^{00}_{\pmb{q}\nu}\right)^{-1}_{\alpha\beta}\\
&=\mathrm{i}\Pi^{00}_{\pmb{q}\nu}\left(\frac{1}{1+U\Pi^{00}_{\pmb{q}\nu}}\delta_{\alpha\beta}+\frac{U\Pi^{00}_{\pmb{q}\nu}}{1-(U\Pi^{00}_{\pmb{q}\nu})^{2}}C_{\alpha\beta}\right)
\end{align*}
新的密度响应函数由下式给出：
\begin{equation*}
\mathrm{i}\Pi^{\mathrm{RPA}^{+}}_{\mathrm{den};\pmb{q}\nu}=\mathrm{i}\Pi^{00}_{\pmb{q}\nu}\left(2\frac{1}{1+U\Pi^{00}_{\pmb{q}\nu}}+2\frac{U\Pi^{00}_{\pmb{q}\nu}}{1-(U\Pi^{00}_{\pmb{q}\nu})^{2}}+2\frac{U\Pi^{00}_{\pmb{q}\nu}}{1-(U\Pi^{00}_{\pmb{q}\nu})^{2}}\right)=2\,\frac{\mathrm{i}\Pi^{00}_{\pmb{q}\nu}}{1-U\Pi^{00}_{\pmb{q}\nu}}
\end{equation*}
新的自旋响应函数则由下式给出：
\begin{equation*}
\mathrm{i}\Pi^{\mathrm{RPA}^{+}}_{\mathrm{spin};\pmb{q}\nu}=\mathrm{i}\Pi^{00}_{\pmb{q}\nu}\left(2\frac{1}{1+U\Pi^{00}_{\pmb{q}\nu}}+2\frac{U\Pi^{00}_{\pmb{q}\nu}}{1-(U\Pi^{00}_{\pmb{q}\nu})^{2}}-2\frac{U\Pi^{00}_{\pmb{q}\nu}}{1-(U\Pi^{00}_{\pmb{q}\nu})^{2}}\right)=\frac{1}{2}\,\frac{\mathrm{i}\Pi^{00}_{\pmb{q}\nu}}{1+U\Pi^{00}_{\pmb{q}\nu}}
\end{equation*}
我们发现 RPA 下的压缩率 $\chi_{c}$ 与自旋磁化率 $\chi_{s}$ 由下式给出：
\begin{align*}
\chi^{\mathrm{RPA}^{+}}_{c}(\pmb{q})&=\frac{\chi^{(0)}_{c}(\pmb{q})}{1+\frac{1}{2}U\chi^{(0)}_{c}(\pmb{q})}\\
\chi^{\mathrm{RPA}^{+}}_{s}(\pmb{q})&=\frac{\chi^{(0)}_{s}(\pmb{q})}{1-2U\chi^{(0)}_{s}(\pmb{q})}
\end{align*}
其中 $\chi^{(0)}_{c}(\pmb{q})=-2\Pi^{00}_{\pmb{q},0}$ 与 $\chi^{(0)}_{s}(\pmb{q})=-\Pi^{00}_{\pmb{q},0}/2$ 是自由电子的压缩率与自旋磁化率。

我们看到，（短程的）排斥相互作用增强自旋磁化率而压制电荷压缩率，而（短程的）吸引相互作用压制自旋磁化率而增强电荷压缩率。对于嵌套费米面与排斥相互作用，当 $T$ 趋于由 $1-U\chi^{(0)}_{s}(\pmb{Q})=0$ 决定的有限临界温度 $T_{c}$ 时，$\chi^{\mathrm{RPA}^{+}}_{s}(\pmb{Q})$ 发散。在 $T_{c}$ 以下，系统自发地产生 SDW，并自发破缺自旋转动对称性与格子的平移对称性。转变温度可以由 $\chi^{(0)}_{s}(\pmb{Q})$ 的低温行为（见式 (5.5.2)）来估计：
\begin{equation*}
T_{c}\sim E_{F}\,\mathrm{e}^{-C_{1}\sqrt{\frac{E_{F}}{U}}}
\end{equation*}
我们看到，对于嵌套费米面，无论相互作用多么弱，排斥相互作用总会产生 SDW。同样，对于嵌套费米面，无论相互作用多么弱，吸引相互作用总会产生 CDW。转变温度为
\begin{equation*}
T_{c}\sim E_{F}\,\mathrm{e}^{-C_{2}\sqrt{\frac{E_{F}}{-U}}}
\end{equation*}

\subsection{自旋密度波态的平均场方法}

\begin{keypoints}
\item 在平均场方法中，我们把费米子算符的对换成它们的基态平均值，从而把一个相互作用的哈密顿量变成一个无相互作用的哈密顿量（平均场哈密顿量）。
\end{keypoints}

我们已经看到，具有嵌套费米面的系统中的排斥相互作用可能引起 SDW 失稳。在本节与接下来的几节中，我们将研究这种系统的基态，并证明基态的确具有 SDW 序。我们将使用两种方法，即平均场方法与变分方法。变分方法在概念上更简单也更正确，它还能计算准粒子之间的相互作用。平均场方法则在数学上更简单。我们先考虑平均场方法。

%FIG{5.20}: {阴影区域是约化布里渊区。它也是半满填充时被填满的费米海。}

为了用平均场方法研究 $T=0$ 的 SDW 态，我们先把哈密顿量改写成更方便的形式。利用恒等式
\begin{equation*}
Un^{2}_{i}=-U(c^{\dagger}_{i}\sigma^{z}c_{i})^{2}+2Un_{i}
\end{equation*}
并把 $2Un_{i}$ 项吸收进化学势项中，我们可以把 Hubbard 模型改写为
\begin{equation*}
H=\sum_{\pmb{k},\alpha}\xi_{\pmb{k}}c^{\dagger}_{\alpha\pmb{k}}c_{\alpha\pmb{k}}-U\sum_{\pmb{i}}(c^{\dagger}_{\pmb{i}}\sigma^{z}c_{\pmb{i}})^{2}
\end{equation*}
现在我们可以形式地得到一个平均场哈密顿量：把自旋 $c^{\dagger}_{i}\sigma^{z}c_{i}$ 之一换成它的平均值 $(-)^{i}M$，并把 $(c^{\dagger}_{i}\sigma^{z}c_{i})^{2}$ 换成 $2M(-)^{i}(c^{\dagger}_{i}\sigma^{z}c_{i})-M^{2}$。我们得到
\begin{equation*}
H^{\mathrm{sub}}_{\mathrm{mean}}=\sum_{k,\alpha}\xi_{k}c^{\dagger}_{\alpha k}c_{\alpha k}-2U\sum_{i}(-)^{i}Mc^{\dagger}_{i}\sigma^{z}c_{i}+N_{\mathrm{site}}UM^{2}
\end{equation*}
在平均场近似下，我们说系统的物理性质由 $H^{\mathrm{sub}}_{\mathrm{mean}}$ 而不是原来的相互作用哈密顿量 (5.5.1) 描述。

平均场哈密顿量可以精确求解。在动量空间中，我们有
\begin{align*}
H^{\mathrm{sub}}_{\mathrm{mean}}&=\sum_{\pmb{k},\alpha}2t(\cos k_{x}+\cos k_{y})c^{\dagger}_{\pmb{k}}c_{\pmb{k}}-B\sum_{\pmb{k}}(-)^{i}c^{\dagger}_{\pmb{k}+\pmb{Q}}\sigma^{z}c_{\pmb{k}}+\text{常数}\\
&=\sum_{\pmb{k}}'\psi^{\dagger}_{\pmb{k}}M_{\pmb{k}}\psi_{\pmb{k}}+\text{常数}
\end{align*}
（译注：第二行中求和号后的 $(-)^{i}$ 为原书排印如此；交错相位因子在过渡到动量空间时已并入下面的 $\psi_{\pmb{k}}$ 与 $M_{\pmb{k}}$ 的定义。）其中 $B=2UM$，$\sum'_{\pmb{k}}$ 表示对约化布里渊区（见图 5.20）求和，$(-)^{i}=(-)^{i_{x}+i_{y}}$，
\begin{equation*}
M_{\pmb{k}}=\begin{pmatrix}\xi_{\pmb{k}} & -B\sigma^{z}\\ -B\sigma^{z} & -\xi_{\pmb{k}}\end{pmatrix},\qquad \xi_{\pmb{k}}=2t(\cos k_{x}+\cos k_{y})\tag{5.5.3}
\end{equation*}
而
\begin{equation*}
\psi_{\alpha\pmb{k}}=\begin{pmatrix}c_{\alpha\pmb{k}}\\ c_{\alpha,\pmb{k}+\pmb{Q}}\end{pmatrix}
\end{equation*}
我们可以通过引入
\begin{align*}
\eta_{\alpha\pmb{k}}&=-v_{\pmb{k}}\sigma^{z}_{\alpha\beta}c_{\beta\pmb{k}}+u_{\pmb{k}}c_{\alpha,\pmb{k}+\pmb{Q}}\\
\lambda_{\alpha\pmb{k}}&=u_{\pmb{k}}c_{\alpha\pmb{k}}+v_{\pmb{k}}\sigma^{z}_{\alpha\beta}c_{\beta,\pmb{k}+\pmb{Q}}
\end{align*}
来对角化上述平均场哈密顿量 $H^{\mathrm{sub}}_{\mathrm{mean}}$，其中
\begin{equation*}
u_{\pmb{k}}=\frac{1}{\sqrt{2}}\sqrt{1-\frac{\xi_{\pmb{k}}}{\sqrt{\xi^{2}_{\pmb{k}}+B^{2}}}}\qquad v_{\pmb{k}}=\mathrm{sgn}(B)\frac{1}{\sqrt{2}}\sqrt{1+\frac{\xi_{\pmb{k}}}{\sqrt{\xi^{2}_{\pmb{k}}+B^{2}}}}
\end{equation*}
由于 $u^{2}_{\pmb{k}}+v^{2}_{\pmb{k}}=1$，$\eta_{\alpha\pmb{k}}$ 与 $\lambda_{\alpha\pmb{k}}$ 满足费米子算符的下列对易关系：
\begin{align*}
\{\eta_{\alpha\pmb{k}},\eta^{\dagger}_{\alpha\pmb{k}'}\}&=\delta_{\pmb{k}\pmb{k}'}\delta_{\alpha\alpha'}\\
\{\lambda_{\alpha\pmb{k}},\lambda^{\dagger}_{\alpha\pmb{k}'}\}&=\delta_{\pmb{k}\pmb{k}'}\delta_{\alpha\alpha'}\\
\{\eta_{\alpha\pmb{k}},\lambda^{\dagger}_{\alpha\pmb{k}'}\}&=0
\end{align*}
用 $\eta_{\alpha\pmb{k}}$ 与 $\lambda_{\alpha\pmb{k}}$ 表示，我们得到
\begin{equation*}
H^{\mathrm{sub}}_{\mathrm{mean}}=\sum'_{\alpha,\pmb{k}}E_{\pmb{k}}\eta^{\dagger}_{\alpha\pmb{k}}\eta_{\alpha\pmb{k}}+\sum'_{\alpha,\pmb{k}}-E_{\pmb{k}}\lambda^{\dagger}_{\alpha\pmb{k}}\lambda_{\alpha\pmb{k}}+N_{\mathrm{site}}UM^{2}
\end{equation*}
而
\begin{equation*}
E_{\pmb{k}}=\sqrt{\xi^{2}_{\pmb{k}}+B^{2}}
\end{equation*}
我们可以看到，平均场基态能量由下式给出：
\begin{equation*}
E_{\mathrm{mean}}=-2\sum'_{\pmb{k}}E_{\pmb{k}}+N_{\mathrm{site}}\frac{B^{2}}{4U}
\end{equation*}
用平均场理论，准粒子激发的谱由 $\pm E_{\pmb{k}}$ 给出。

然而，上述平均场结果包含一个未知参量 $M$——交错自旋的平均值。为了确定 $M$，我们使平均场基态能量 $E_{\mathrm{mean}}$ 取极小。我们得到
\begin{equation*}
\frac{1}{4}-UN^{-1}_{\mathrm{site}}\sum'_{\pmb{k}}\frac{1}{E_{\pmb{k}}}=0.\tag{5.5.4}
\end{equation*}
这个方程称为能隙方程。通过施加自洽关系
\begin{equation*}
M=(-)^{i}\langle \Phi_{\mathrm{mean}}|c^{\dagger}_{i}\sigma^{z}c_{i}|\Phi_{\mathrm{mean}}\rangle
\end{equation*}
我们也可以得到同一个能隙方程，其中 $|\Phi_{\mathrm{mean}}\rangle$ 是 $H^{\mathrm{sub}}_{\mathrm{mean}}$ 的基态。也就是说，由平均场基态得到的平均自旋应当等于我们在得到平均场哈密顿量时所假设的平均自旋。

如果能隙方程有 $M\neq 0$ 的解，那么我们这个相互作用系统的基态就自发破缺自旋转动对称性，并发展出反铁磁序。由于 $\int\frac{\mathrm{d}^{2}k}{(2\pi)^{2}}\frac{1}{E_{k}}$ 在 $B\to 0$ 时趋于无穷，我们发现只要 $U>0$，式 (5.5.4) 总有 $M\neq 0$ 的解。于是，正 $U$ 的 Hubbard 模型总会发展出 SDW，无论 $U$ 多么小。这与上一节得到的结果一致。

我们想指出，只要 $\xi_{\pmb{k}}=-\xi_{\pmb{k}+\pmb{Q}}$，上述结果就仍然成立。在这种情形下，费米面（即 $\xi_{\pmb{k}}=0$ 处）上不同的点由 $\pmb{Q}$ 相连；换言之，费米面满足嵌套条件。对于更一般的 $\xi_{\pmb{k}}$，即费米面不满足嵌套条件的情形，准粒子的色散由下式给出：
\begin{equation*}
\pm\sqrt{\frac{1}{4}(\xi_{\pmb{k}}-\xi_{\pmb{k}+\pmb{Q}})^{2}+B^{2}}+\frac{1}{2}(\xi_{\pmb{k}}+\xi_{\pmb{k}+\pmb{Q}})
\end{equation*}
在这种情形下，能隙方程 (5.5.4) 不再适用。事实表明，$U$ 较小时并不会形成 SDW。

虽然平均场方法（以及下面讨论的变分方法）正确地重现了排斥型 Hubbard 模型的反铁磁基态，但由此得到的激发谱却是不正确的。这是因为反铁磁基态自发破缺了自旋转动对称性，基态之上应当存在无能隙的自旋波激发。自旋波激发可以通过自旋关联 $\mathrm{i}S^{ab}(i-j)=\left\langle (c^{\dagger}_{i}\sigma^{a}c_{i})(c^{\dagger}_{j}\sigma^{b}c_{j})\right\rangle$ 来探测。在平均场理论中，我们有 $S^{ab}\sim\delta_{ab}G^{2}_{\mathrm{mean}}$，其中 $G_{\mathrm{mean}}$ 是由 $H_{\mathrm{mean}}$ 确定的电子格林函数。于是，只有当 $|\omega|>2\Delta$ 时 $\mathrm{Im}S^{ab}_{\pmb{k}\omega}$ 才非零。然而，如果我们超越平均场理论，用 RPA 以及图 5.18 与图 5.19 中的梯形图来计算 $S^{ab}$，那么 $\mathrm{Im}S^{ab}_{\pmb{k}\omega}$ 将一直到 $\omega=0$ 都非零，这表明存在无能隙激发。此外，$S^{ab}_{\pmb{k}\omega}$ 还有一个孤立极点，由它我们可以得到自旋波激发的色散。

\subsection{自旋密度波态的变分方法——艰难的途径}

\begin{keypoints}
\item 在变分方法中，我们用一个无相互作用哈密顿量的基态作为尝试波函数，去近似一个相互作用哈密顿量的基态。
\end{keypoints}

在变分方法中，我们关注的是基态波函数，而不是哈密顿量。我们已经看到，具有在位排斥相互作用的 Hubbard 模型有 SDW 失稳。为了用变分方法理解基态的性质，我们索性为基态猜测一个尝试波函数。一个可能的选择是像我们在 Hartree–Fock 近似中所做的那样，使用 $H_{0}=\sum_{\pmb{k},\alpha}\tilde{\xi}_{\pmb{k}}c^{\dagger}_{\alpha\pmb{k}}c_{\alpha\pmb{k}}$ 的基态 $|\Psi_{0}\rangle$ 作为尝试波函数。然而，$|\Psi_{0}\rangle$ 并不破缺自旋与平移对称性。SDW 失稳的存在提示我们，应当能够找到一个能量更低的不同的尝试波函数。从对称性的考虑出发，我们愿意使用
\begin{equation*}
H_{B}=\sum_{\pmb{k},\alpha}2\tilde{t}(\cos k_{x}+\cos k_{y})c^{\dagger}_{\alpha\pmb{k}}c_{\alpha\pmb{k}}-B\sum_{i}(-)^{i}c^{\dagger}_{i}\sigma^{z}c_{i}
\end{equation*}
的基态 $|\Psi_{B}\rangle$ 作为尝试波函数。我们并不显式写出多体尝试波函数，而是像在平均场方法中那样，先对角化"变分"哈密顿量 $H_{B}$。$H_{B}$ 的基态满足 $\lambda^{\dagger}_{\alpha\pmb{k}}|\Psi_{B}\rangle=\eta_{\alpha\pmb{k}}|\Psi_{B}\rangle=0$。这一代数关系确定了我们的尝试波函数。

注意，哈密顿量 $H_{B}$ 只是用来得到尝试波函数的一个算符。把 $H_{B}$ 缩放一个常数不会改变尝试波函数。因此我们可以令 $\tilde{t}=t$，从而 $\tilde{\xi}_{\pmb{k}}=\xi_{\pmb{k}}$。此时 $H_{B}$ 与平均场哈密顿量 $H^{\mathrm{sub}}_{\mathrm{mean}}$ 具有相同的形式。这里 $B$ 是一个改变波函数形状的变分参量，我们稍后将改变 $B$ 使能量极小。

为了计算尝试态的平均能量，我们注意到，对于尝试波函数 $|\Psi_{B}\rangle$，我们有 $\langle\eta^{\dagger}_{\alpha\pmb{k}}\eta_{\alpha\pmb{k}}\rangle=0$ 与 $\langle\lambda^{\dagger}_{\alpha\pmb{k}}\lambda_{\alpha\pmb{k}}\rangle=1$。由
\begin{equation*}
c_{\alpha\pmb{k}}=u_{\pmb{k}}\lambda_{\alpha\pmb{k}}-v_{\pmb{k}}\sigma^{z}_{\alpha\beta}\eta_{\beta\pmb{k}}\qquad c_{\alpha,\pmb{k}+\pmb{Q}}=v_{\pmb{k}}\sigma^{z}_{\alpha\beta}\lambda_{\beta\pmb{k}}+u_{\pmb{k}}\eta_{\alpha\pmb{k}}
\end{equation*}
我们发现
\begin{align*}
\left\langle c^{\dagger}_{\alpha\pmb{k}}c_{\beta\pmb{k}}\right\rangle&=u^{2}_{\pmb{k}}\delta_{\alpha\beta} & \left\langle c^{\dagger}_{\alpha,\pmb{k}+\pmb{Q}}c_{\beta,\pmb{k}+\pmb{Q}}\right\rangle&=v^{2}_{\pmb{k}}\delta_{\alpha\beta}\\
\left\langle c^{\dagger}_{\alpha\pmb{k}}c_{\beta,\pmb{k}+\pmb{Q}}\right\rangle&=u_{\pmb{k}}v_{\pmb{k}}\sigma^{z}_{\alpha\beta} & \left\langle c^{\dagger}_{\alpha,\pmb{k}+\pmb{Q}}c_{\beta\pmb{k}}\right\rangle&=u_{\pmb{k}}v_{\pmb{k}}\sigma^{z}_{\alpha\beta}
\end{align*}
在实空间中，我们有
\begin{equation*}
\left\langle c^{\dagger}_{\alpha i}c_{\beta i}\right\rangle=\frac{1}{2}\delta_{\alpha\beta}+2(-)^{i}\sigma^{z}_{\alpha\beta}N^{-1}_{\mathrm{site}}\sum'_{\pmb{k}}u_{\pmb{k}}v_{\pmb{k}}=\frac{1}{2}\delta_{\alpha\beta}+(-)^{i}\sigma^{z}_{\alpha\beta}N^{-1}_{\mathrm{site}}\sum'_{\pmb{k}}\frac{B}{E_{\pmb{k}}}
\end{equation*}
正如所料，我们的尝试态 $|\Psi_{B}\rangle$ 具有反铁磁序：
\begin{equation*}
\langle S_{z,i}\rangle=\frac{1}{2}\left\langle c^{\dagger}_{i}\sigma^{z}c_{i}\right\rangle=(-)^{i}M,\qquad M=N^{-1}_{\mathrm{site}}\sum'_{\pmb{k}}\frac{B}{E_{\pmb{k}}}
\end{equation*}

我们尝试态的平均能量由下式给出（见 5.5.4.1 节）：
\begin{equation*}
\langle\Psi_{B}|H|\Psi_{B}\rangle=N_{\mathrm{site}}\left(-2N^{-1}_{\mathrm{site}}\sum'_{\pmb{k}}\frac{\xi^{2}_{\pmb{k}}}{E_{\pmb{k}}}-U\left(N^{-1}_{\mathrm{site}}\sum'_{\pmb{k}}\frac{B}{E_{\pmb{k}}}\right)^{2}+\text{常数}\right)
\end{equation*}
$B$ 的值通过使 $\langle\Psi_{B}|H|\Psi_{B}\rangle$ 极小而得到。我们发现这样的 $B$ 满足能隙方程（见 5.5.4.2 节）：
\begin{equation*}
1-UN^{-1}_{\mathrm{site}}\sum'_{\pmb{k}}\frac{1}{\sqrt{\xi^{2}_{\pmb{k}}+B^{2}}}=0\tag{5.5.5}
\end{equation*}

我们看到式 (5.5.4) 与式 (5.5.5) 并不一致。这说明了这样一点：平均场方法在领头阶就可能给出不同的结果。这是因为不同的平均场方法对应于不同的展开，而这些展开一般具有不同的领头项。

变分方法也让我们能够研究激发的性质。激发的尝试波函数是在尝试基态 $|\Psi_{B}\rangle$ 的基础上由 $\eta^{\dagger}_{\alpha\pmb{k}}$ 与 $\lambda_{\alpha\pmb{k}}$ 产生的。对于激发态，我们有 $\left\langle\eta^{\dagger}_{\alpha\pmb{k}}\eta_{\alpha\pmb{k}}\right\rangle=n^{+}_{\alpha\pmb{k}}$ 与 $\left\langle\lambda^{\dagger}_{\alpha\pmb{k}}\lambda_{\alpha\pmb{k}}\right\rangle=n^{-}_{\alpha\pmb{k}}$，其中 $n^{+}_{\alpha\pmb{k}}\neq 0$ 且 $n^{-}_{\alpha\pmb{k}}\neq 1$。把 $\langle c^{\dagger}c\rangle$ 用 $n^{+}_{\alpha\pmb{k}}$ 与 $n^{-}_{\alpha\pmb{k}}$ 表出之后，我们得到尝试激发态的平均能量（见 5.5.4.3 节）如下：
\begin{equation*}
\langle\Psi^{\mathrm{exc}}_{B}|H|\Psi^{\mathrm{exc}}_{B}\rangle=E_{\mathrm{ground}}+\sum'_{\pmb{k},\alpha}E_{\pmb{k}}(\delta n^{+}_{\alpha\pmb{k}}-\delta n^{-}_{\alpha\pmb{k}})+U\delta N+O(\delta n^{2})
\end{equation*}
其中 $n^{-}_{\alpha\pmb{k}}=1+\delta n^{-}_{\alpha\pmb{k}}$ 与 $n^{+}_{\alpha\pmb{k}}=\delta n^{+}_{\alpha\pmb{k}}$ 描述围绕基态的小涨落。我们看到准粒子激发具有能谱 $\pm E_{\pmb{k}}$。注意，准粒子激发存在有限的能隙 $\Delta=\min E_{\pmb{k}}=|B|$。准粒子之间的相互作用可以通过 $(\delta n^{\pm}_{\alpha\pmb{k}})^{2}$ 诸项方便地包括进来。

\subsection{评注：一些计算细节}

\subsubsection{平均能量的计算}

我们如下计算尝试态 $|\Psi_{B}\rangle$ 的平均能量：
\begin{align*}
&\langle\Psi_{B}|H|\Psi_{B}\rangle\\
&\quad=2\sum'_{\pmb{k}}(u^{2}_{\pmb{k}}-v^{2}_{\pmb{k}})\xi_{\pmb{k}}+\frac{U}{2}\sum_{i}\left(\left\langle c^{\dagger}_{\alpha i}c_{\alpha i}\right\rangle\left\langle c^{\dagger}_{\beta i}c_{\beta i}\right\rangle+\left\langle c^{\dagger}_{\alpha i}c_{\beta i}\right\rangle\left\langle c_{\alpha i}c^{\dagger}_{\beta i}\right\rangle\right)\\
&\quad=2\sum'_{\pmb{k}}(u^{2}_{\pmb{k}}-v^{2}_{\pmb{k}})\xi_{\pmb{k}}+\frac{U}{2}\sum_{i}\left(1+\mathrm{Tr}\big(\frac{1}{2}+(-)^{i}M\sigma^{z}\big)\big(\frac{1}{2}-(-)^{i}M\sigma^{z}\big)\right)\\
&\quad=N_{\mathrm{site}}\left(-UM^{2}+\frac{3U}{4}+2N^{-1}_{\mathrm{site}}\sum'_{\pmb{k}}-\frac{\xi_{\pmb{k}}}{E_{\pmb{k}}}\xi_{\pmb{k}}\right)
\end{align*}

\subsubsection{能隙方程的计算}

为了得到使平均能量取极小的 $B$ 值，我们计算 $\partial\langle\Psi_{B}|H|\Psi_{B}\rangle/\partial B$，得到
\begin{align*}
&2N^{-1}_{\mathrm{site}}\sum'_{\pmb{k}}-\xi^{2}_{\pmb{k}}\delta E^{-1}_{\pmb{k}}-2U\left(N^{-1}_{\mathrm{site}}\sum'_{\pmb{k}}\frac{B}{E_{\pmb{k}}}\right)\left(N^{-1}_{\mathrm{site}}\sum'_{\pmb{k}}\frac{\delta B}{E_{\pmb{k}}}-\frac{B^{2}\delta B}{E^{3}_{\pmb{k}}}\right)\\
&\quad=2N^{-1}_{\mathrm{site}}\sum'_{\pmb{k}}-\xi^{2}_{\pmb{k}}\delta E^{-1}_{\pmb{k}}-2U\left(N^{-1}_{\mathrm{site}}\sum'_{\pmb{k}}\frac{B}{E_{\pmb{k}}}\right)\left(N^{-1}_{\mathrm{site}}\sum'_{\pmb{k}}\frac{\xi^{2}_{\pmb{k}}\delta B}{E^{3}_{\pmb{k}}}\right)\\
&\quad=\left(2N^{-1}_{\mathrm{site}}\sum'_{\pmb{k}}-\xi^{2}_{\pmb{k}}\delta E^{-1}_{\pmb{k}}\right)-2U\left(N^{-1}_{\mathrm{site}}\sum'_{\pmb{k}}\frac{B}{E_{\pmb{k}}}\right)\left(N^{-1}_{\mathrm{site}}\sum'_{\pmb{k}}-\xi^{2}_{\pmb{k}}B^{-1}\delta E^{-1}_{\pmb{k}}\right)
\end{align*}
上式为零的要求给出能隙方程 (5.5.5)。

\subsubsection{激发态平均能量的计算}

由关系
\begin{align*}
\left\langle c^{\dagger}_{\alpha\pmb{k}}c_{\beta\pmb{k}}\right\rangle&=u^{2}_{\pmb{k}}n^{-}_{\alpha\pmb{k}}\delta_{\alpha\beta}+v^{2}_{\pmb{k}}n^{+}_{\alpha\pmb{k}}\delta_{\alpha\beta}, & \left\langle c^{\dagger}_{\alpha,\pmb{k}+\pmb{Q}}c_{\beta,\pmb{k}+\pmb{Q}}\right\rangle&=v^{2}_{\pmb{k}}n^{-}_{\alpha\pmb{k}}\delta_{\alpha\beta}+u^{2}_{\pmb{k}}n^{+}_{\alpha\pmb{k}}\delta_{\alpha\beta},\\
\left\langle c^{\dagger}_{\alpha\pmb{k}}c_{\beta,\pmb{k}+\pmb{Q}}\right\rangle&=u_{\pmb{k}}v_{\pmb{k}}(n^{-}_{\alpha\pmb{k}}-n^{+}_{\alpha\pmb{k}})\sigma^{z}_{\alpha\beta}, & \left\langle c^{\dagger}_{\alpha,\pmb{k}+\pmb{Q}}c_{\beta\pmb{k}}\right\rangle&=u_{\pmb{k}}v_{\pmb{k}}(n^{-}_{\alpha\pmb{k}}-n^{+}_{\alpha\pmb{k}})\sigma^{z}_{\alpha\beta},
\end{align*}
我们发现，在实空间中，
\begin{equation*}
\left\langle c^{\dagger}_{\alpha i}c_{\beta i}\right\rangle=\delta_{\alpha\beta}N^{-1}_{\mathrm{site}}\sum'_{\pmb{k}}(n^{-}_{\alpha\pmb{k}}+n^{+}_{\alpha\pmb{k}})+2(-)^{i}\sigma^{z}_{\alpha\beta}N^{-1}_{\mathrm{site}}\sum'_{\pmb{k}}u_{\pmb{k}}v_{\pmb{k}}(n^{-}_{\alpha\pmb{k}}-n^{+}_{\alpha\pmb{k}})
\end{equation*}
% ---------------------------------------------------------------------------
% 块边界备注（装配说明，非译文）：
% 1. 起点：书页 226（p241）顶部 "charges in a metal…" 是书页 225（p240）末句
%    "However, we know that two" 的续半句。ch5_c 正文已止于"然而，我们知道，金属中的"，
%    其续文即本文件第 2 行 %JOIN% 之后的"两个电荷由于来自其他电子的屏蔽，……"，
%    与 ch5_c 文件头 %--C-TAIL-- 注释中的建议译文一致。
% 2. 终点：本块止于书页 238（p253）末（5.5.4.3 的实空间关系式）。书页 239（p254）
%    顶部 "This allows us to calculate" 起的推导与式 (5.5.6) 属 ch5_e（书页 239–249，
%    见工作日志块划分）。按 assemble.py 约定（X-TAIL 区只放上一块的收尾），
%    本文件不设 %--C-TAIL-- 收尾内容区，也勿把 p254 内容译入本块；
%    请 ch5_e 以其正文首行 %JOIN%（"这使我们可以计算"）续接本块末式。
% 3. %JOIN% 必须是本文件第 2 行（首行注释之后不能有别的内容行或注释行），
%    否则 assemble.py 的缝合检测会失效——此为本次块的重要装配约束。
% 4. 蓝本问题备注：蓝本把半满填充写作 ⟨⟨n_i⟩⟩=1，PNG 为 ⟨n_i⟩=1，已按 PNG 校正；
%    蓝本两处 T_c 公式（C_1 与 E_F/U、C_2 与 E_F/(−U)）经放大核对与 PNG 一致；
%    蓝本 5.5.2 节动量空间平均场哈密顿量中的 (−)^i 与 PNG 一致（原书排印如此），
%    已照录并加译注。
